\begin{itemframe}{طراحی الگوریتم‌های تقریبی}
\itm
در این بخش دو تکنیک قدرتمند برای طراحی الگوریتم‌های تقریبی را بررسی می‌کنیم:
تصادفی‌سازی و برنامه‌ریزی خطی.
\itm
ابتدا با یک الگوریتم تصادفی برای نسخه‌ی بهینه‌سازی مسئله‌ی ارضاء پذیری 3-CNF آغاز می‌کنیم
و سپس نشان می‌دهیم چگونه می‌توان یک الگوریتم تقریبی برای نسخه‌ی وزن‌دار مسئله‌ی پوشش رأسی طراحی کرد که مبتنی بر برنامه‌ریزی خطی باشد.
\end{itemframe}


\begin{itemframe-s}{طراحی الگوریتم‌های تقریبی}{طراحی به وسیله الگوریتم‌های تصادفی}
\itm
همان‌طور که برخی الگوریتم‌های تصادفی پاسخ دقیق محاسبه می‌کنند، برخی دیگر پاسخ‌های تقریبی تولید می‌کنند.
\itm
می‌گوییم یک الگوریتم تصادفی برای مسئله‌ای دارای ضریب تقریبی برابر با
$\alpha(n)$
است، اگر:
$$
\max \left\{ \frac{C}{C^{\ast}}, \frac{C^{\ast}}{C} \right\} \leq \alpha(n).
$$
n
اندازه‌ی ورودی و C هزینه‌ی مورد انتظار است.
\itm
به بیان دیگر، الگوریتم تقریب تصادفی مشابه الگوریتم تقریب قطعی است، با این تفاوت که ضریب تقریب بر اساس هزینه‌ی مورد انتظار بیان می‌شود.
\end{itemframe-s}


\begin{itemframe-s}{طراحی الگوریتم‌های تقریبی}{طراحی به وسیله الگوریتم‌های تصادفی}
\itm
پیش‌تر با مسئله‌ی ارضاء‌پذیری 3-CNF آشنا شدیم. برای ارضاء‌پذیری عبارت، باید انتسابی از متغیرها وجود داشته باشد که هر بند
\fn{clause}
 به مقدار ۱ ارزیابی شود.
\itm
اگر نمونه‌ای ارضاء‌پذیر نباشد، شاید بخواهید بدانید که تا چه حد به ارضاء‌پذیر بودن نزدیک است؛ یعنی انتسابی از متغیرها بیابیم که بیشترین تعداد بندها را ارضاء کند. این مسئله‌ی بیشینه‌سازی را ارضاء‌پذیری
MAX-3-CNF
می‌نامیم.

\end{itemframe-s}


\begin{itemframe-s}{طراحی الگوریتم‌های تقریبی}{طراحی به وسیله الگوریتم‌های تصادفی}
\itm
شاید شگفت‌زده شوید که اگر هر متغیر با احتمال 1/2 به ۱ و با احتمال 1/2 به ۰ مقداردهی شود، الگوریتمی با ضریب تقریب
$\frac{8}{7}$
 به‌دست می‌آید. در ادامه دلیل آن را خواهیم دید.
\itm
در اینجا بر خلاف تعریف ارضاء‌پذیری 3-CNF فرض می‌کنیم که هیچ بندی شامل یک متغیر و نقیض آن نباشد.

\itm
فرض کنید n متغیر با نام‌های
$x_1, x_2, \dots, x_n$
 و m بند داشته باشیم و هر متغیر به‌طور مستقل با احتمال
$\frac{1}{2}$
 برابر ۱ و با احتمال
$\frac{1}{2}$
 برابر ۰ قرار گیرد.
حال برای هر
$i = 1, 2, \dots, m$
، متغیر تصادفی نشانگر زیر را تعریف کنید:
$$Y_i = I\{\text{ ارضاء شود } i \text{ بند }\}$$
\end{itemframe-s}


\begin{itemframe-s}{طراحی الگوریتم‌های تقریبی}{طراحی به وسیله الگوریتم‌های تصادفی}
\itm
پس
$Y_i = 1$
است زمانی که حداقل یکی از گزاره‌نماهای
\fn{literal}
بند i برابر ۱ باشد.
از آن‌جا که هیچ گزاره‌نمایی بیش از یک بار در یک بند ظاهر نمی‌شود، و همچنین فرض کرده‌ایم که هیچ متغیر و نقیض آن در یک بند نیامده باشند، مقادیر سه گزاره‌نما در هر بند مستقل‌اند.
\itm
یک بند تنها زمانی ارضاء نمی‌شود که هر سه گزاره‌نمای آن برابر صفر باشند. بنابراین:
$$
\Pr\{\text{ ارضاء نشود} i \text{ بنده }\} = (1/2)^3 = 1/8
$$
در نتیجه:
$$
\Pr \{\text{ ارضاء شود } i \text{ بند }\} = 1 - 1/8 = 7/8
$$
\end{itemframe-s}


\begin{itemframe-s}{طراحی الگوریتم‌های تقریبی}{طراحی به وسیله الگوریتم‌های تصادفی}
\itm
بر اساس ویژگی‌های متغییرهای تصادفی نشانگر داریم:
$$
E[Y_i] = 7/8
$$
حال اگر Y تعداد کل بندهای ارضاشده باشد داریم:
$$
Y = Y_1 + Y_2 + \dots + Y_m
$$

\end{itemframe-s}


\begin{itemframe-s}{طراحی الگوریتم‌های تقریبی}{طراحی به وسیله الگوریتم‌های تصادفی}
\itm
پس باید امید ریاضی متغییر  Y را محاسبه کنیم:
\begin{align*}
E[Y] &= E\left[\sum_{i=1}^m Y_i\right] \\[2pt]
     &= \sum_{i=1}^m E[Y_i]  \\[2pt]
     &= \sum_{i=1}^m 7/8 \\[2pt]
     &= 7m/8.
\end{align*}

\end{itemframe-s}


\begin{itemframe-s}{طراحی الگوریتم‌های تقریبی}{طراحی به وسیله الگوریتم‌های تصادفی}

\itm
از آن‌جا که m یک کران بالا برای تعداد بندهای ارضاء‌پذیر است، ضریب تقریب در بیشترین حالت برابر خواهد بود با:
$$
\frac{m}{7m/8} = 8/7
$$


\end{itemframe-s}


\begin{itemframe-s}{طراحی الگوریتم‌های تقریبی}{طراحی به وسیله برنامه ریزی خطی}
\itm
در مسئله‌ی «پوشش رأسی با وزن کمینه‌»
\fn{minimum-weight vertex-cover problem}
هر رأس مانند v وزنی مثبت به نام $w(v)$ دارد.
وزن یک پوشش رأسی مانند
$V'$
به‌صورت زیر تعریف می‌شود:
$$
w(V') = \sum_{v \in V'} w(v)
$$
هدف، یافتن یک پوشش رأسی با وزن کمینه است.
\itm
الگوریتم تقریبی که برای حالت بدون وزن دیدیم، در اینجا کارآمد نیست.
زیرا برای حالت وزن‌دار جواب حاصل می‌تواند فاصله‌ی زیادی از جواب بهینه داشته باشد.
\end{itemframe-s}


\begin{itemframe-s}{طراحی الگوریتم‌های تقریبی}{طراحی به وسیله برنامه ریزی خطی}
\itm
در عوض با استفاده از یک برنامه‌ریزی خطی، یک کران پایین برای وزن پوشش رأسی محاسبه می‌کنیم.
سپس این جواب را گرد کرده و از آن برای به‌دست‌آوردن یک پوشش رأسی استفاده می‌کنیم.
\itm
ابتدا برای هر رأس $v \in V$ یک متغیر $x(v)$ تعریف می‌کنیم و الزام می‌کنیم که $x(v)$ برابر با $0$ یا $1$ باشد.
رأس $v$ تنها در صورتی در پوشش رأسی وجود دارد که
$$x(v) = 1$$
باشد. بنابراین، قید مربوط به هر یال $(u, v)$ این است که دست‌کم یکی از $u$ و $v$ در پوشش رأسی قرار داشته باشند، که می‌توان آن را به صورت زیر نوشت:
$$x(u) + x(v) \geq 1$$
\end{itemframe-s}


\begin{itemframe-s}{طراحی الگوریتم‌های تقریبی}{طراحی به وسیله برنامه ریزی خطی}
\itm
این دیدگاه منجر به یک «برنامه‌ریزی صحیح 0-1»
\fn{0-1 integer program}
 برای یافتن پوشش رأسی کمینه می‌شود:
\begin{align*}
\text{minimize } & \sum_{v \in V} w(v) x(v)\\
\text{ to subject }& \\[3pt]
&x(u) + x(v) \geq 1 \quad \text{ each for} (u, v) \in E\\[3pt]
&x(v) \in \{0, 1\} \quad \text{ each for} v \in V
\end{align*}
\end{itemframe-s}


\begin{itemframe-s}{طراحی الگوریتم‌های تقریبی}{طراحی به وسیله برنامه ریزی خطی}
\itm
یک حالت خاص از صورت بندی بالا را در نظر بگیرید که تمام وزن‌ها $w(v) = 1$ باشند.  این صورت‌بندی معادل نسخه‌ی بهینه ‌سازی مسئله‌ی پوشش رأسی است. می‌دانیم که پوشش رأسی ان‌پی کامل و نسخه بهینه سازی آن ان‌پی سخت است.
\itm
مشخص است که این مسئله از نسخه بهینه‌سازی مسئله پوشش رأسی سخت‌تر است پس حداقل ان‌پی سخت است.

\end{itemframe-s}


\begin{itemframe-s}{طراحی الگوریتم‌های تقریبی}{طراحی به وسیله برنامه ریزی خطی}
\itm
اما ما به دنبال حل دقیق این مسئله نیستیم پس بیایید قید
$
x(v) \in \{0, 1\}
$
را با
$$
0 \leq x(v) \leq 1
$$

جایگزین کنیم. در نتیجه، برنامه‌ریزی خطی زیر به‌دست می‌آید:
\begin{align*}
\text{ minimize } &\sum_{v \in V} w(v) x(v)\\
\text{ to subject}& \\
&x(u) + x(v) \geq 1 \quad \text{ each for} (u, v) \in E \\
&x(v) \leq 1 \quad \text{ each for} v \in V \\
&x(v) \geq 0 \quad \text{ each for} v \in V\\
\end{align*}

\end{itemframe-s}


\begin{itemframe-s}{طراحی الگوریتم‌های تقریبی}{طراحی به وسیله برنامه ریزی خطی}
\itm
به این کار «ساده‌سازی»
\fn{relaxation}
 برنامه‌ریزی خطی گفته می‌شود.
هر جواب مجاز برای برنامه‌ی اصلی، یک جوابی مجاز برای نسخه ساده شده است است.
\itm
بنابراین، مقدار یک جواب بهینه در برنامه‌ریزی خطی ساده شده، یک کران پایین برای مقدار جواب بهینه در برنامه‌ی اصلی است.
\itm
در نتیجه مقدار یک جواب بهینه در برنامه‌ریزی خطی ساده شده، یک کران پایین برای وزن بهینه در مسئله‌ی «پوشش رأسی با وزن کمینه»‌ است.
\end{itemframe-s}


\begin{itemframe-s}{طراحی الگوریتم‌های تقریبی}{طراحی به وسیله برنامه ریزی خطی}
\itm
روش
\texttt{APPROX-MIN-WEIGHT-VC}
 با استفاده از یک جواب برای نسخه ساده‌شده برنامه‌ریزی خطی آغاز می‌کند و از آن برای ساختن یک جواب تقریبی برای مسئله‌ی پوشش رأسی کمینه استفاده می‌کند:

\begin{algorithm}[H]\alglr
\caption{APPROX-MIN-WEIGHT-VC}
\begin{algorithmic}[1]
\Func{APPROX-MIN-WEIGHT-VC}{$G, w$}
\State $C \gets \emptyset$
\State compute $\bar{x}$, an optimal solution to the linear-programming relaxation
\For{each vertex $v \in V$}
    \If{$\bar{x}(v) \geq \frac{1}{2}$}
        \State $C \gets C \cup \{v\}$
    \EndIf
\EndFor
\State \Return $C$
\end{algorithmic}
\end{algorithm}

\end{itemframe-s}


\begin{itemframe-s}{طراحی الگوریتم‌های تقریبی}{طراحی به وسیله برنامه ریزی خطی}
\itm
این روش به شکل زیر عمل می‌کند: خط 2 ساده‌شده برنامه‌ریزی خطی را صورت‌بندی کرده و آن را حل می‌کند. یک جواب بهینه که حاصل حل برنامه‌ریزی خطی است، به هر رأس مانند $v$ یک مقدار
$\bar{x}(v)$
بین ۰ و ۱ نسبت می‌دهد.
\itm
سپس این مقدار برای انتخاب رئوسی که باید به پوشش رأسی افزوده شوند به‌کار می‌رود: رأس $v$ تنها در صورتی در $C$ قرار می‌گیرد که
$$
\bar{x}(v) \geq 1/2
$$
باشد.

در واقع، روش هر متغیر را به $0$ یا $1$ گرد می‌کند تا جوابی برای برنامه‌ی خطی اصلی به‌دست آید که جوابی برای مسئله پوشش رأسی است.
\end{itemframe-s}


\begin{itemframe-s}{طراحی الگوریتم‌های تقریبی}{طراحی به وسیله برنامه ریزی خطی}
\itm
الگوریتم
\texttt{APPROX-MIN-WEIGHT-VC}
 یک الگوریتم تقریبی با ضریب 2 و در زمان چندجمله‌ای برای مسئله‌ی پوشش رأسی با وزن کمینه است.
\itm
از آن‌جا که الگوریتمی در زمان چندجمله‌ای برای حل برنامه‌ریزی خطی (در خط 2) وجود دارد و حلقه‌ی for نیز در زمان چندجمله‌ای اجرا می‌شود، این الگوریتم چندجمله‌ای است.
\end{itemframe-s}


\begin{itemframe-s}{طراحی الگوریتم‌های تقریبی}{طراحی به وسیله برنامه ریزی خطی}
\itm
اکنون باید نشان دهیم که این الگوریتم یک الگوریتم تقریبی با ضریب 2 است. فرض کنید
$C^*$
یک جواب بهینه برای مسئله‌ی پوشش رأسی کمینه باشد، و
$z^*$
مقدار یک جواب بهینه برای نسخه ساده‌شده آن باشد.
\itm
 از آن‌جا که یک پوشش رأسی بهینه جوابی مجاز برای برنامهٔ ساده‌شده است، داریم:
$$
z^* \leq w(C^*)
$$
\end{itemframe-s}


\begin{itemframe-s}{طراحی الگوریتم‌های تقریبی}{طراحی به وسیله برنامه ریزی خطی}
\itm
اما چگونه گرد کردن مقادیر کسری $\bar{x}(v)$، پوشش رأسی است؟
\itm
برای دیدن این موضوع، یالی مانند
$(u, v)$
را در نظر بگیرید. در برنامه خطی ساده‌شده قید کردیم:
$$
x(u) + x(v) \geq 1
$$
\itm
پس حداقل یکی از $\bar{x}(u)$ یا $\bar{x}(v)$ باید بزرگ‌تر یا مساوی با $1/2$ باشد. بنابراین، حداقل یکی از $u$ یا $v$ در پوشش رأسی قرار می‌گیرد و همه‌ی یال‌ها پوشش داده می‌شوند.
\end{itemframe-s}


\begin{itemframe-s}{طراحی الگوریتم‌های تقریبی}{طراحی به وسیله برنامه ریزی خطی}
\itm
حال می‌‌خواهیم ثابت کنیم:
$$
w(C) \leq 2 w(C^*)
$$
برای اینکار ابتدا باید ثابت کنیم:
$$
w(C) \leq 2 z^*
$$
\end{itemframe-s}


\begin{itemframe-s}{طراحی الگوریتم‌های تقریبی}{طراحی به وسیله برنامه ریزی خطی}
\itm
بیایید از وزن پوشش شروع کنیم:
\begin{align*}
z^* &= \sum_{v \in V} w(v) \bar{x}(v)\\[3pt]
&\geq \sum_{v \in V : \bar{x}(v) \geq 1/2} w(v) \bar{x}(v)\\[3pt]
&\geq \sum_{v \in V : \bar{x}(v) \geq 1/2} w(v) \cdot \tfrac{1}{2}\\[3pt]
&= \tfrac{1}{2} \sum_{v \in C} w(v)\\[3pt]
&= \tfrac{1}{2} w(C)
\end{align*}

\end{itemframe-s}


\begin{itemframe-s}{طراحی الگوریتم‌های تقریبی}{طراحی به وسیله برنامه ریزی خطی}
\itm
پیش‌تر دیدیم که
$$
z^* \leq 2 w(C^*)
$$
پس داریم:
$$
w(C) \leq 2 z^* \leq 2 w(C^*)
$$
\itm
بنابراین
$$
w(C) \leq 2 w(C^*)
$$
\itm
در نتیجه، الگوریتم
\texttt{APPROX-MIN-WEIGHT-VC}
الگوریتم دارای ضریب تقریب 2 است.
\end{itemframe-s}